\chapter[Industrial Defect Calculus]{Industrial Defect Calculus for High-Dimensional Nonlinear Filtering}
\label{ch:bf-highdim-candidate-synthesis}

\section{Why The Synthesis Must Start With Defects}
\label{sec:bf-hd-synthesis-purpose}

The preceding chapters explain high-dimensional nonlinear filtering methods.
An industrial macro-finance program needs a harsher object: a calculus of
failure.  A method is useful for a global-bank or central-bank model only when
the failure mode is explicit, the mitigation has a mathematical contract, and
the cost variable that can kill the method is visible before implementation.
Academic novelty is not the target here.  The target is a defensible
composition of known methods under stated assumptions.

This chapter therefore treats every candidate as a conditional component.  A
Gaussian scaffold is useful if it exposes local curvature and covariance
failure.  Sparse-grid or high-degree cubature is useful if it diagnoses local
nonlinearity without pretending to be a global filter.  Tensor methods are
useful if economic structure keeps ranks stable and preserves probability or
covariance semantics.  Transport maps are useful if they reduce geometry while
leaving a checkable target.  HMC is useful only downstream of a same-scalar
likelihood contract.

The reader-facing synthesis map is this: start with the exact conditional law
and likelihood normalizers from Chapter~\ref{ch:bf-highdim-nonlinear-foundations};
replace them only by named approximations; attach one diagnostic to each
object being approximated; and promote a method only after the previous object
exports the input required by the next.  The chapter's exact target is not a
new algorithm.  It is an auditable decision procedure for composing existing
methods without losing the scalar likelihood, probability-law, support,
Jacobian, rank, or covariance semantics needed by industrial inference.

\paragraph{Running quadratic-observation cell.}
The scalar cell
\[
  x_t=\rho x_{t-1}+\eta_t,\qquad y_t=x_t^2+\epsilon_t
\]
has now served four roles.  Chapter~\ref{ch:bf-highdim-nonlinear-foundations}
used it to show that the exact filtering law can have two lobes.
Chapter~\ref{ch:bf-highdim-gaussian-quadrature} used it to show why Gaussian
projection can preserve moments while losing shape.
Chapter~\ref{ch:bf-highdim-particle-transport-tensor} used it to separate
particles, transports, and tensor representations.
Chapter~\ref{ch:bf-highdim-hmc-research} used it to insist that HMC needs a
scalar likelihood and a same-scalar gradient.  Here the same cell is not
evidence for a macro-finance model.  It is the small mechanism that tells a
mixed numerical reader what the architecture is trying not to lose.

\paragraph{Imports from the preceding chapters.}
\begingroup
\scriptsize
\setlength{\tabcolsep}{2pt}
\begin{longtable}{@{}p{0.16\linewidth}p{0.20\linewidth}p{0.21\linewidth}p{0.18\linewidth}p{0.16\linewidth}@{}}
\toprule
Source chapter & Exported object & Exported defect & Exported diagnostic/veto
& Consumed here as \\
\midrule
\endhead
Ch.~\ref{ch:bf-highdim-nonlinear-foundations} & conditional law, normalizer,
likelihood gradient & wrong target, missing normalizer, value-gradient
mismatch & finite \(Z_t\), mass, positivity, same-scalar check & target and
likelihood gates \\
Ch.~\ref{ch:bf-highdim-gaussian-quadrature} & Gaussian moment pair and
quadrature rule & projection bias, exactness without accuracy, point explosion
& innovation, PSD, active dimension, comparator residual, approximate-score
parity & scaffold, local curvature diagnostic, and approximate HMC target gate \\
Ch.~\ref{ch:bf-highdim-particle-transport-tensor} & empirical measure,
transport map, TT/TN representation & collapse, support failure, rank/semantic
failure, PSD loss, nonsmooth adaptive branches & ESS, support/logdet, rank,
mass, positivity, factor gate, HMC admissibility label & conditional
non-Gaussian tools \\
Ch.~\ref{ch:bf-highdim-hmc-research} & scalar potential, gradient, transformed
target & same-scalar mismatch, divergences, invalid transform & scalar parity,
finite support/Jacobian, energy diagnostics & downstream sampler gate \\
\bottomrule
\end{longtable}
\endgroup

\section[Particle Collapse]{Defect 1: Particle Collapse Is Exponential Without Structure}
\label{sec:bf-hd-defect-particles}

The bootstrap and SIR filters in Chapter~\ref{ch:bf-highdim-particle-transport-tensor}
are honest baselines because their collapse mechanism is visible.  Let
$W_i=\exp(L_i)$ be unnormalized weights and
$\bar W_i=W_i/\sum_{j=1}^N W_j$.  The effective sample size is
\begin{equation}
  \operatorname{ESS}=\frac{1}{\sum_{i=1}^N \bar W_i^2}.
  \label{eq:bf-hd-synth-ess}
\end{equation}
Suppose the log-weight decomposes across conditionally independent observation
blocks,
\begin{equation}
  L_i=\sum_{k=1}^{d_{\mathrm{eff}}} \ell_{ik},\qquad
  \E[\ell_{ik}]=\mu,\qquad \Var(\ell_{ik})=\sigma_\ell^2,
  \label{eq:bf-hd-synth-logweight}
\end{equation}
with weak dependence sufficient for
\(\Var(L_i)\asymp d_{\mathrm{eff}}\sigma_\ell^2\).  This is the informal
calculus behind the high-dimensional particle-filter warnings of
\citet{bengtsson2008curse} and \citet{snyder2008obstacles}.

\begin{proposition}[Log-weight variance defect]
\label{prop:bf-hd-particle-collapse-calculus}
If the log-weight variance grows like
\(\Var(L_i)=v_d\asymp d_{\mathrm{eff}}\sigma_\ell^2\), then a particle filter
with proposal independent of the current observation needs \(N\) to grow on the
scale of \(\exp\{c v_d\}\), for a model-dependent constant \(c>0\), to keep
more than a small number of normalized weights from dominating.
\end{proposition}

\begin{proof}[Derivation]
For a lognormal heuristic, \(\E[W]=\exp(\mu_d+v_d/2)\) and
\(\E[W^2]=\exp(2\mu_d+2v_d)\), so
\[
  \frac{\E[W^2]}{\E[W]^2}=\exp(v_d).
\]
The usual importance-sampling variance factor is therefore exponential in
\(v_d\).  Since the denominator of \eqref{eq:bf-hd-synth-ess} is driven by the
largest normalized weights, stable ESS requires the sample size to offset this
exponential dispersion.  The cited papers give sharper asymptotic statements;
the monograph uses this calculation to expose the industrial defect: without
locality, guided proposals, or likelihood tempering, dimension enters through
log-likelihood variance, not through a benign constant.
\end{proof}

The mitigation is not "use more particles" unless the cost model permits
exponential growth.  The viable mitigations are structural:
\begin{enumerate}[label=(\arabic*)]
  \item reduce \(d_{\mathrm{eff}}\) by block-local assimilation;
  \item reduce \(\sigma_\ell^2\) by guided or transport proposals;
  \item use pseudo-marginal or particle-MCMC logic only when the likelihood
    estimator and variance contract are explicit
    \citep{andrieu2009pseudo,andrieu2010particle};
  \item veto by ESS, maximum weight, unique ancestors, and likelihood-variance
    diagnostics before comparing runtime.
\end{enumerate}

A concrete Gaussian observation cell makes the exponential defect visible.  If
\(Y=0\), \(X\sim\calN(0,I_d)\), and the likelihood is
\(g(y\mid x)\propto \exp\{-\|x\|^2/(2R)\}\), then
\[
  L=\log g(0\mid X)=C-\frac{1}{2R}\sum_{k=1}^d X_k^2,\qquad
  \Var(L)=\frac{d}{2R^2}.
\]
The lognormal variance factor above is therefore \(\exp\{d/(2R^2)\}\).  This
example is deliberately simple; its role is to show why a large macro-finance
observation panel cannot be rescued by a particle count that scales only
linearly with \(d\).

\section{Defect 2: Gaussian And Cubature Exactness Is Not Posterior Accuracy}
\label{sec:bf-hd-defect-gaussian}

Gaussian filters fail by projection.  If \(X\sim\calN(m,P)\) and a quadrature
rule is exact for a polynomial class \(\mathcal P_p\), then
\begin{equation}
  \E[\varphi(X)]-\sum_j w_j\varphi(m+L\xi_j)
  =
  \E[r_p(X)]-\sum_j w_j r_p(m+L\xi_j),
  \qquad
  \varphi=\Pi_p\varphi+r_p .
  \label{eq:bf-hd-synth-quadrature-remainder}
\end{equation}
This identity is harmless only when the non-polynomial remainder \(r_p\) is
small on the predictive support.  In macro-finance likelihoods, kinks from
occasionally binding constraints, effective lower bounds, regime switches, or
measurement transformations can make \(r_p\) the main object.

\begin{proposition}[Local diagnostic role of high-order rules]
\label{prop:bf-hd-local-cubature-diagnostic}
Sparse-grid and high-degree cubature rules are industrially defensible first as
local curvature diagnostics.  They become global filter components only when a
block decomposition bounds the active dimension and when the higher-order
moment change alters a decision-relevant likelihood, score, or innovation
diagnostic.
\end{proposition}

\begin{proof}[Derivation]
The CKF, UKF, high-degree CKF, and SGQF sources
\citep{julier1997new,arasaratnam2009cubature,jia2012sparsegrid,jia2013highdegree}
construct deterministic rules for Gaussian-weighted moment integrals.  The
filter still projects the updated law back to a Gaussian moment pair.  Thus a
high-order rule reduces the quadrature part of
\eqref{eq:bf-hd-synth-quadrature-remainder}; it does not remove Gaussian
projection error.  If a block has dimension \(b\), the point count and dense
linear algebra are controlled by \(b\), not by the full state dimension \(n\).
Without such a block, improved polynomial exactness can simply buy an
unaffordable global rule.
\end{proof}

The mitigation is a two-level test: use the Gaussian scaffold as the default
stabilizer, then run high-order or sparse-grid diagnostics only on blocks whose
innovation residuals, score sensitivity, or Hessian contractions indicate
nonlinear curvature.  \citet{singh2018adaptive} is useful here as an adaptive
sparse-grid competitor, but its empirical/adaptive construction is not a
general guarantee for nonsmooth industrial likelihoods.

Projection error is not a subtle tail event.  The symmetric mixture
\[
  p_a(x)=\frac12\phi(x;-a,s^2)+\frac12\phi(x;a,s^2)
\]
has mean \(0\) and variance \(a^2+s^2\), so a Gaussian projection records only
\(\calN(0,a^2+s^2)\).  As \(a/s\) grows, the posterior becomes visibly
bimodal while the projected law remains unimodal.  No polynomial cubature rule
around the projected Gaussian changes that semantic loss; it can only improve
the moment integrals used before the projection.

In the running quadratic-observation cell, this is exactly what the observation
\(y_t=x_t^2+\epsilon_t\) does when \(y_t>0\) and \(R\) is small: it makes the
state sign hard to identify while making \(|x_t|\) informative.  A cubature
rule can compute Gaussian-weighted moments of \(x^2\) accurately and still
passes a single Gaussian lobe to the next method.  The synthesis therefore
uses Chapter~\ref{ch:bf-highdim-gaussian-quadrature} in a precise way: it
computes Gaussian projection residuals, innovation covariances, and
same-scalar likelihood gradients on selected blocks, then refuses to interpret
those quantities as a global posterior representation when the block posterior
is visibly multimodal.

\section[Tensor Compression Defects]{Defect 3: Tensor Compression Can Destroy The Object Being Compressed}
\label{sec:bf-hd-defect-tensor}

Tensor trains replace \(M^d\) grid storage by \(O(dMr^2)\) storage when the
TT rank is at most \(r\) \citep{oseledets2011tt}.  The condition "when \(r\)
is small" is not decoration; it is the entire theorem-to-implementation bridge.
For a filtering density \(p(x)\), tensor compression can fail in four distinct
ways:
\begin{equation}
  \text{rank growth},\qquad
  \int \widehat p(x)\,dx \ne 1,\qquad
  \widehat p(x)<0,\qquad
  \widehat\ell_T(\theta)\not\approx \ell_T(\theta).
  \label{eq:bf-hd-synth-tt-defects}
\end{equation}
For covariance compression, the analogous defect is indefiniteness:
\begin{equation}
  \widehat P=\operatorname{round}_{\mathrm{TN}}(P),\qquad
  \lambda_{\min}(\widehat P)<0.
  \label{eq:bf-hd-synth-tn-indefinite}
\end{equation}
Even an entrywise small covariance perturbation can leave the PSD cone.  For
\[
  P=\begin{pmatrix}1&1\\1&1\end{pmatrix},\qquad
  \widehat P_\epsilon=\begin{pmatrix}1&1+\epsilon\\1+\epsilon&1\end{pmatrix},
\]
the eigenvalues of \(\widehat P_\epsilon\) are \(2+\epsilon\) and
\(-\epsilon\).  Thus a tiny unconstrained rounding error can make a covariance
unusable as a Gaussian filter object.

\begin{proposition}[Tensor viability requires stable structure]
\label{prop:bf-hd-tensor-viability}
A tensor method is a viable industrial filtering component only when the
economic model induces a stable low-rank structure for the represented object
and when the compressed object passes the semantic checks appropriate to that
object: mass and positivity for densities, square-root or PSD checks for
covariances, and likelihood/score agreement for posterior targets.
\end{proposition}

\begin{proof}[Derivation]
The TT storage count \(O(dMr^2)\) follows from storing \(d\) cores of mode size
\(M\) and adjacent rank bounded by \(r\).  If \(r\) grows exponentially with
coupling or likelihood interaction, the apparent dimensional saving disappears.
For a density, rounding error is not only an approximation norm: it may change
sign or mass and hence leave the probability simplex.  For a covariance,
rounding entries or tensor cores does not preserve the cone
\(\{P:P\succeq0\}\).  A factor representation \(P=SS^\top\) preserves PSD
algebraically, which is the lesson behind square-root Kalman filtering and the
tensor-network square-root warning of \citet{menzen2024tnsrkf}.  Finally, HMC
and parameter inference use a scalar target; a small tensor residual is not
sufficient unless the induced likelihood and score errors are controlled.
\end{proof}

This proposition explains the synthesis.  The Li--Wang--Yau--Zhang and
Fox--Dolgov--Morrison--Molteno routes are density/PDE pilots
\citep{li2019ttfiltering,fox2021grid}; Zhao--Cui is the discrete-time
TT/KR state-parameter learning route \citep{zhao2024ttsequential}; Meng
preprints mark provisional PR-DMZ and correlated-noise routes
\citep{meng2025regularity,meng2026correlated}; Batselier--Chen--Wong and
Menzen--Kok--Batselier support lifted-linear and square-root covariance
cautions \citep{batselier2016tnkf,menzen2024tnsrkf}.  None of these sources
alone validates a general nonlinear tensor filter for BayesFilter.

\section{Defect 4: Transports Can Improve Geometry While Breaking Auditability}
\label{sec:bf-hd-defect-transport}

Let \(x=T(z)\) be a differentiable bijection.  The transformed density is
\begin{equation}
  \pi_Z(z)=\pi_X(T(z))|\det DT(z)|.
  \label{eq:bf-hd-synth-transport-density}
\end{equation}
This identity is the mathematical boundary between a useful transport and a
dangerous one.  A learned or fitted map can reduce dependence and curvature,
but if its support, monotonicity, Jacobian, or target correction is wrong, it
changes the inference problem.

\begin{proposition}[Transport usefulness requires target auditability]
\label{prop:bf-hd-transport-auditability}
A transport map is industrially useful only if it reduces a measured geometric
defect while preserving a checkable target identity: either an explicit
importance/proposal correction, a deterministic analysis-map contract, or a
Jacobian-corrected transformed posterior.
\end{proposition}

\begin{proof}[Derivation]
Equation~\eqref{eq:bf-hd-synth-transport-density} follows by change of
variables.  For proposal filtering, the same calculation gives the proposal
density and the density-ratio correction.  For MCMC, it gives the transformed
potential used by transport-map MCMC and NeuTra-style HMC
\citep{parno2018transport,hoffman2019neutra}.  For ensemble filtering,
\citet{reich2013nonparametric} supports the narrower deterministic transform
from forecast to equally weighted analysis particles, while
\citet{spantini2022coupling} supports stochastic and deterministic map filters
with localization.  In all cases, the map is a computational object, not a
certificate.  The certificate is the target identity plus residual diagnostics.
\end{proof}

The mitigation is to report map residuals, finite inverse/Jacobian checks,
localization sensitivity, ESS or analysis-spread diagnostics, and reference
moments on controlled cells.  A lower transport loss or smoother geometry is
not enough if the corrected target cannot be reconstructed.

The support condition is equally concrete.  If the target density
\(\gamma(x)\) is positive on both half-lines but a proposal or transported
reference density \(q(x)\) is zero for \(x<0\), then the correction
\(\gamma(x)/q(x)\) is infinite on a set of positive target mass.  No downstream
HMC, ESS, or map residual can repair a proposal whose support excludes part of
the target.

For the running quadratic-observation cell, a proposal or map that covers only
the positive lobe near \(+\sqrt y\) may look locally excellent.  It is still
invalid for the two-lobed target unless the excluded negative lobe has zero
target mass or an explicit approximation decision has been made.  This is why
transport auditability is a support and correction question before it is an
efficiency question.

\section{Defect 5: HMC Is Downstream Of The Filter Target}
\label{sec:bf-hd-defect-hmc}

For parameters \(\theta=\tau(q)\), HMC uses a scalar potential
\begin{equation}
  U(q)=-\ell_T(\tau(q))-\log\pi(\tau(q))
       -\log|\det D\tau(q)|.
  \label{eq:bf-hd-synth-hmc-target}
\end{equation}
If the filter supplies \(\widehat\ell_T\), the sampler targets
\(\widehat\pi(\theta\mid y_{1:T})\), not the exact posterior, unless a
pseudo-marginal or exact correction is part of the algorithm
\citep{andrieu2009pseudo,andrieu2010particle}.
The analytical exact-gradient contract is developed in
Chapter~\ref{ch:bf-highdim-hmc-research}; an approximate-filter gradient
supports only the declared approximate scalar.  It is not an exact posterior
gradient or exact state-space likelihood gradient unless an exact correction
proof is supplied.

\begin{proposition}[HMC is downstream of a same-scalar likelihood contract]
\label{prop:bf-hd-hmc-downstream}
HMC, RMHMC, NUTS, NeuTra, and transport-preconditioned HMC should be compared
only after the filter defines the same finite scalar value used for gradients,
accept/reject decisions, diagnostics, and posterior interpretation.
\end{proposition}

\begin{proof}[Derivation]
The detailed-balance argument for HMC in \citet{neal2011mcmc} uses a
Hamiltonian whose gradient generates the numerical proposal and whose value
appears in the Metropolis correction.  If the gradient is not the derivative of
the reported scalar, the proposal is no longer a discretization of the accepted
Hamiltonian.  If the scalar is approximate, the chain may still be valid for
the approximate target, but the posterior claim has changed.  RMHMC
\citep{girolami2011riemann}, NUTS \citep{hoffman2014nuts}, and NeuTra
\citep{hoffman2019neutra} add geometry or adaptation; they do not remove this
same-scalar requirement.
\end{proof}

The mitigation ladder is therefore: finite value and gradient checks, scalar
parity, fixed-HMC baseline, divergence and energy diagnostics
\citep{betancourt2017conceptual}, posterior/reference checks on controlled
models, then transport or metric acceleration.  Speed enters after validity.
A minimal scalar-parity diagnostic is
\[
  \Delta_j(q;h)=
  \frac{U_{\mathrm{reported}}(q+h e_j)-U_{\mathrm{reported}}(q-h e_j)}{2h}
  -[\nabla U_{\mathrm{used}}(q)]_j .
\]
If \(\max_j|\Delta_j(q;h)|\) is not small on typical parameter states, the
sampler is not using the gradient of the scalar it reports.

\section{Numerical Defects And Mitigation Equations}
\label{sec:bf-hd-numerical-defects}

\small
\begin{longtable}{@{}p{0.16\linewidth}p{0.30\linewidth}p{0.30\linewidth}p{0.15\linewidth}@{}}
\toprule
Defect & Mathematical symptom & Mitigation contract & Veto diagnostic \\
\midrule
\endhead
Log-weight underflow &
\(\log\sum_i e^{L_i}=m+\log\sum_i e^{L_i-m}\), \(m=\max_iL_i\) &
Use log-sum-exp for arithmetic stability; still reject if
\(\operatorname{ESS}/N\) collapses &
ESS, max weight, unique ancestors \\
Ill-conditioned innovation &
\(S=HPH^\top+R\), \(\kappa(S)\gg1\) &
Solve by factorization; regularize only with declared likelihood change &
failed Cholesky, large residual \( \|SK-C\| \) \\
PSD loss &
\(\lambda_{\min}(\widehat P)<0\) &
Use Joseph/square-root/factor updates; check reconstructed PSD &
negative eigenvalue, failed factor solve \\
TT density drift &
\(\int\widehat p\ne1\) or \(\min_x\widehat p(x)<0\) &
Track mass, positivity, truncation error; use squared/nonnegative
representations where source-supported &
mass error, negative cells, rank spike \\
Support mismatch &
\(q(x)=0\) where \(\gamma(x)>0\) &
Require support coverage and finite density ratio &
infinite log weight, failed inverse or Jacobian \\
Scalar mismatch &
\(\nabla U_{\mathrm{used}}\ne\nabla U_{\mathrm{reported}}\) &
Differentiate the same scalar used for accept/reject &
finite-difference mismatch, divergences \\
\bottomrule
\end{longtable}
\normalsize

These mitigations are not evidence of posterior accuracy.  They are necessary
conditions for an industrial method to remain mathematically auditable.

\section{Performance Models And Killing Variables}
\label{sec:bf-hd-performance-models}

\small
\begin{longtable}{@{}p{0.15\linewidth}p{0.30\linewidth}p{0.27\linewidth}p{0.18\linewidth}@{}}
\toprule
Component & Nominal cost model & Killing variable & Mitigation \\
\midrule
\endhead
Dense Gaussian filter &
\(O(n^3)\) factorization, \(O(n^2)\) memory &
full state dimension \(n\) &
block, factor, or sparse solves \\
Sigma point or CKF &
\(O(N_q C_h)+O(n^3)\), CKF \(N_q=2n\) &
state dimension and observation cost \(C_h\) &
block diagnostics, square-root factors \\
Sparse grid or high degree &
\(O(N_{\mathrm{sg}}(b,\ell)C_h)\) for block dimension \(b\), level \(\ell\) &
effective block dimension and smoothness &
adaptive/local blocks only \\
Particle methods &
\(O(NTC_f)\) model calls; stable \(N\) may scale like \(\exp(c d_{\mathrm{eff}})\) &
effective observation dimension &
guided, localized, or transport proposals \\
TT density or PDE &
\(O(dMr^2)\) storage plus rank-dependent operator cost &
rank \(r\), grid mode \(M\), boundary error &
rank diagnostics, local subproblems \\
Transport maps &
ensemble/map-fit cost plus logdet/Jacobian cost &
map dimension, polynomial/order, ensemble size &
triangular/local maps and residual checks \\
HMC through the filter &
\(O(N_{\mathrm{grad}}\, C_{\ell,\nabla\ell})\) &
gradient-through-filter cost and invalid regions &
same-scalar checks, fixed baseline, then acceleration \\
\bottomrule
\end{longtable}
\normalsize

The performance rule is simple: compare enhanced methods only after the defect
vetoes pass.  A fast invalid target or an inexpensive collapsed particle cloud
is not an industrial solution.

For sparse grids the symbol \(N_{\mathrm{sg}}(b,\ell)\) hides the practical
danger.  In the source-local SGQF formulation of \citet{jia2012sparsegrid},
the point count is controlled relative to a tensor product only for a declared
level and active dimension; it is not a dimension-free filter.  Because the
original Smolyak, Stroud, and Genz foundations remain scoped as source risks in
this lane, the industrial mitigation is stated narrowly: spend sparse-grid
evaluations only where a block diagnostic has made \(b\), level, smoothness,
and omitted-interaction risk visible.

\section{Synthesis Contract Architecture}
\label{sec:bf-hd-synthesis-architecture}

The preceding chapters pass explicit defects into the architecture:
Chapter~\ref{ch:bf-highdim-nonlinear-foundations} passes the exact-target and
normalizer contract; Chapter~\ref{ch:bf-highdim-gaussian-quadrature} passes
projection and quadrature defects;
Chapter~\ref{ch:bf-highdim-particle-transport-tensor} passes collapse,
correction, rank, positivity, and PSD defects; and
Chapter~\ref{ch:bf-highdim-hmc-research} passes the same-scalar and
Jacobian-corrected target defects.  The architecture is therefore not a ranked
list of fashionable methods.  It is a sequence of vetoes.

Let \(D_j\) denote a method-specific defect diagnostic and let
\(\tau_j\) denote the threshold or qualitative gate declared before a run.  A
component \(M_{j+1}\) is admissible only if the preceding component exports
enough information to check the assumptions of \(M_{j+1}\):
\[
  D_j\le \tau_j
  \quad\Longrightarrow\quad
  \text{the input object required by }M_{j+1}\text{ is well defined.}
  \label{eq:bf-hd-contract-handoff}
\]
This implication is an engineering contract, not a new convergence theorem.
It says exactly what must be audited before a method handoff can be defended.

\begin{proposition}[Sparse-grid promotion requires bounded active dimension]
\label{prop:bf-hd-sparse-grid-promotion}
A sparse-grid or high-degree cubature rule should remain a local diagnostic
unless the active block dimension, smoothness class, point budget, and
omitted-interaction residual are all declared and checked.
\end{proposition}

\begin{proof}[Derivation]
The checked sparse-grid and high-degree cubature filtering papers construct
Gaussian-weighted deterministic integration rules for stated levels, degrees,
and assumptions.  Their cost is controlled by the active dimension entering the
rule, not by a slogan that sparse grids beat dimension.  If a nonlinear
observation block depends effectively on \(b\) variables, the relevant cost is
\(N_{\rm sg}(b,\ell)C_h\) or the corresponding high-degree point count.  If
the same map couples all \(n\) variables, \(b=n\) and the local diagnostic
becomes a global expensive filter.  The promotion gate is therefore a bound on
\(b\), a smoothness check, and a residual showing that omitted interactions do
not drive a decision-relevant scalar.
\end{proof}

\begin{proposition}[Performance claims require validity gates first]
\label{prop:bf-hd-performance-after-veto}
Runtime, memory, ESS per gradient, rank, and point count can compare methods
only after the target, semantic, and numerical veto diagnostics have passed.
\end{proposition}

\begin{proof}[Derivation]
Let \(V\) be a validity event: finite scalar likelihood, value-gradient parity,
noncollapsed weights, support coverage, PSD covariance, or density mass and
positivity, depending on the component.  A performance statistic \(C\) is
interpretable for the intended inference problem only on \(V\).  On \(V^c\),
the computation may be fast because it solves the wrong problem, collapses to a
single particle, rounds a covariance outside the PSD cone, or samples an
approximate target that is not being reported.  Thus the admissible comparison
is conditional: compare \(C\mid V\), and treat \(V^c\) as a veto rather than a
slow datapoint.
\end{proof}

\section{Controlled Industrial-Style Worked Example}
\label{sec:bf-hd-industrial-worked-example}

The scalar running cell now scales into a controlled macro-finance stress cell
by replacing one latent scalar with a block state and replacing \(x^2\) by a
block nonlinear observation.  The analogy is useful but limited: the scalar
cell teaches the defects; it does not validate their severity in a real central
bank or global-bank model.

Consider a stylized macro-finance state
\[
  x_t=(m_t,f_t,v_t,u_t),\qquad
  m_t\in\R^{b_m},\quad f_t\in\R^{b_f},\quad
  v_t\in\R,\quad u_t\in\R^{b_u},
\]
where \(m_t\) are macro factors, \(f_t\) are financial-stress factors,
\(v_t\) is a latent volatility state, and \(u_t\) are measurement disturbances
or survey wedges.  Suppose a Gaussian scaffold supplies a forecast
\(\calN(a_t,P_t^-)\) with block factors, and a stress observation has nonlinear
mean
\[
  h(x_t)=\beta^\top m_t + f_t^\top A f_t + \exp(v_t/2)u_{t,1},
  \qquad
  y_t=h(x_t)+\varepsilon_t,\quad
  \varepsilon_t\sim\calN(0,R).
  \label{eq:bf-hd-worked-stress-cell}
\]
This is not a NAWM validation, a client model, or production evidence.  It is a
controlled cell for deciding which defects are active.

The Gaussian scaffold first computes \(\bar y\), \(S\), and \(C\) as in
Chapter~\ref{ch:bf-highdim-gaussian-quadrature}.  If the block factorization
fails, \(S\) is ill conditioned, or the posterior covariance loses PSD, the
cell has already failed the cheapest audit.  If it passes, the pass only says
that the local moment machinery is numerically coherent.

A particle diagnostic asks whether an unstructured proposal is even plausible.
For \(d_{\rm eff}\) approximately independent stress indicators with a
Gaussian cell like the example in Section~\ref{sec:bf-hd-defect-particles},
\[
  \Var(L)\approx \frac{d_{\rm eff}}{2R^2},
  \qquad
  \frac{\E[W^2]}{\E[W]^2}\approx
  \exp\!\left\{\frac{d_{\rm eff}}{2R^2}\right\}.
  \label{eq:bf-hd-worked-pf-factor}
\]
If this factor is large, the architecture does not ask for a heroic global
particle count.  It requires a block-local likelihood, guided proposal,
transport correction, tempering contract, or rejection of the particle route
for this cell.

A sparse-grid diagnostic is then local to \(z_t=(f_t,v_t,u_{t,1})\), whose
dimension is \(b=b_f+2\), not the full model dimension.  The diagnostic is
allowed only if the map in \eqref{eq:bf-hd-worked-stress-cell} is smooth on the
forecast support and if \(N_{\rm sg}(b,\ell)C_h\) fits the budget.  The output
is not a global sparse-grid filter; it is a comparison of Gaussian and
high-order moment contributions to \(\widehat Z_t\), the score, or the
innovation residual.

The tensor route is scoped by rank and semantics.  If a TT pilot represents a
local density over \(z_t\), the admissibility row is
\[
  r_{\max}\le r_{\rm cap},\qquad
  \left|\int \widehat p(z)\,dz-1\right|\le \epsilon_{\rm mass},\qquad
  \min_z \widehat p(z)\ge -\epsilon_+,
  \label{eq:bf-hd-worked-tt-gate}
\]
plus a likelihood or moment agreement check.  A rank spike or negative density
cell blocks promotion even if storage is small.

The transport/HMC route is downstream.  A triangular map may be useful if its
support, monotonicity, and Jacobian checks pass, but the transformed target
must still be
\[
  U_z(z)=-\widehat\ell_T(\tau(T(z)))-\log\pi(\tau(T(z)))
        -\log|\det D\tau(T(z))|-\log|\det DT(z)|.
  \label{eq:bf-hd-worked-hmc-target}
\]
The scalar-parity diagnostic from Chapter~\ref{ch:bf-highdim-hmc-research} is
then applied to this same scalar.  Only after these gates pass may the program
compare fixed HMC, transport-preconditioned HMC, or another sampler by ESS per
gradient or wall time.

The worked example therefore demonstrates defect scoping.  It does not choose
a validated method.  It says which defects would force a global bank or central
bank implementation team to stop, localize, correct, or fall back before
calling a method industrially meaningful.

\paragraph{Reader checkpoint.}
What should now be clear: the synthesis is not ranking methods by fashion or
novelty.  Object preserved: a chain of auditable handoffs from law and
normalizer to approximation, diagnostic, and sampler target.  Approximation
enters at: each method-specific representation and each promotion gate.
Exported to implementation: promote only after the target, semantic, numerical,
and performance vetoes for the relevant object have passed.

\section{Industrial Synthesis Propositions}
\label{sec:bf-hd-industrial-propositions}

\begin{proposition}[Block scaffold before exotic methods]
\label{prop:bf-hd-block-scaffold-first}
A large nonlinear macro-finance filter should first expose a block-local
Gaussian scaffold before adding particles, tensors, transports, or HMC.  The
scaffold is not trusted because it is accurate; it is trusted because its
innovation, covariance, derivative, and factorization failures are cheap to
observe.
\end{proposition}

\begin{proof}[Derivation]
The Gaussian update computes a moment projection using linear solves and
covariance factors.  These objects produce residuals, condition numbers,
innovation sequences, and PSD checks at cost far below global particle or HMC
experiments.  If this scaffold fails on a controlled block, a more complicated
method has no reliable local baseline.  If it passes, the pass is only a
diagnostic: multimodality and nonlocal dependence can still be missed.
\end{proof}

\begin{proposition}[Useful synthesis need not be novel]
\label{prop:bf-hd-useful-not-novel}
A non-novel composition can be industrially superior to a novel standalone
method when each component removes a specific defect and exports a diagnostic
that the next component can test.
\end{proposition}

\begin{proof}[Derivation]
Let component \(A\) reduce a defect measure \(D_A\) while exporting residual
\(R_A\), and component \(B\) require \(R_A\le\tau_A\) before its own target is
well defined.  The composition is meaningful only through the implication
\[
  R_A\le\tau_A \quad\Longrightarrow\quad
  \text{the assumptions of component } B \text{ are approximately testable}.
\]
For example, block Gaussian diagnostics can identify curvature blocks for
sparse-grid checks; a TT density pilot can supply conditional-density structure
for a KR transport; a transport can precondition HMC only after the
Jacobian-corrected scalar is audited.  None of these compositions is claimed as
a new theorem; the value is that every handoff has a defect variable and a
veto.
\end{proof}

\begin{algorithm}
\caption{Defect-first high-dimensional filtering synthesis}
\label{alg:bf-hd-defect-synthesis}
\begin{algorithmic}[1]
\Require Model partition, source-supported method assumptions, controlled
  references, defect veto thresholds.
\State Run block Gaussian scaffold; record innovation, factor, derivative, and
  conditioning residuals.
\State Apply high-order or sparse-grid diagnostics only to blocks with local
  curvature or innovation defects.
\State If particle collapse is predicted by log-weight variance, add guided or
  transport proposals and retain explicit correction.
\State If density/PDE structure is plausible, run TT/FTT pilots only on blocks
  with stable rank, mass, positivity, and boundary diagnostics.
\State If covariance compression is needed, prefer square-root or factor
  representations and veto on PSD/factor failure.
\State Define the scalar likelihood/posterior target and verify value-gradient
  consistency.
\State Run fixed HMC baseline; add RMHMC, NeuTra, transport maps, or TT/KR
  acceleration only after target diagnostics pass.
\State Promote no component whose performance win occurs before validity
  gates pass.
\end{algorithmic}
\end{algorithm}

\section{Literature Coverage And Source Risks}
\label{sec:bf-hd-source-risks}

The defect calculus relies on locally inspected source families recorded in the
P1R--P2R ledgers.  The strongest source-local support covers particle collapse,
Gaussian/cubature competitors, TT/TN filtering routes, transport-map filtering,
and HMC/transport preconditioning.  Remaining source risks are explicit:
Savostyanov-specific maxvol quasioptimality, Stroud's book, Smolyak and Genz
originals, Knothe's original rearrangement, some local-particle-filter
literature, and broader numerical-analysis references such as Higham-style
finite-precision PSD analysis have not been source-local checked in this lane.
They may motivate future literature expansion, but they do not support
theorem-level claims here.

\section{Formal And Editorial Boundaries}
\label{sec:bf-hd-formal-boundaries}

The propositions in this chapter are monograph derivations and industrial
contracts, not full paper-length proofs.  MathDevMCP diagnostics for the P4 and
P5 labels are recorded in the corresponding MathDevMCP audit ledgers; they
locate the relevant obligations and certify only narrow algebraic checks, not
broad propositions.  Any remaining PDF layout warnings are editorial issues,
not mathematical support.  A panel should therefore read the chapter as a
source-grounded synthesis with explicit diagnostics, not as a machine-verified
theorem file or a completed implementation validation.

\section{What Is Not Concluded}
\label{sec:bf-hd-synth-nonclaims}

This chapter does not conclude that high-dimensional nonlinear filtering is
solved for BayesFilter, that any method is NAWM-ready, that HMC converges, that
tensor or transport methods are validated, that posterior accuracy is
established, that GPU/XLA is ready, or that any method should become a
production default.  It concludes only that an industrial candidate must pass a
defect calculus before speed, elegance, or novelty is allowed to matter.
